
\documentclass{article}
\usepackage{colortbl}
\usepackage{makecell}
\usepackage{multirow}
\usepackage{supertabular}

\begin{document}

\newcounter{utterance}

\centering \large Interaction Transcript for game `eqbench', experiment `eq\_bench', episode 60 with reflexion{-}Qwen3.5{-}27B.
\vspace{24pt}

{ \footnotesize  \setcounter{utterance}{1}
\setlength{\tabcolsep}{0pt}
\begin{supertabular}{c@{$\;$}|p{.15\linewidth}@{}p{.15\linewidth}p{.15\linewidth}p{.15\linewidth}p{.15\linewidth}p{.15\linewidth}}
    \# & \multicolumn{2}{c}{Player} && \multicolumn{2}{c}{Game Master} \\
    \hline

    \theutterance \stepcounter{utterance}  
    & & & \multicolumn{4}{p{0.6\linewidth}}{
        \cellcolor[rgb]{0.9,0.9,0.9}{
            \makecell[{{p{\linewidth}}}]{
                \texttt{\tiny{[P1$\langle$GM]}}
                \texttt{You have attempted this same task before and did not succeed. These are the plans you wrote after those attempts.} \\
\texttt{Trial \#1: New plan: My previous attempt failed because I misjudged the primary emotional driver of the final interaction. I scored "Surprise" and "Shame" very high, but the dialogue reveals that Tyler has already exposed Claudia's secret (the inheritance) and she admits she didn't know he knew. This implies the surprise was immediate and has likely subsided into the realization of her failure. The core of the scene is Tyler accusing her of abandonment and selfishness ("You abandoned me," "You never bothered to find out"), which directly targets her guilt. Therefore, "Shame" should be the dominant emotion, but "Surprise" should be lower as the shock of discovery has passed. "Anger" is low because she is defensive and pleading, not aggressive. "Anticipation" is minimal as the conversation is ending on a negative note. Next time, I will prioritize the emotion that aligns with the character's admission of fault and the accuser's specific grievances, ensuring the highest score reflects the weight of the accusation (Shame) rather than the initial shock (Surprise). I will adjust the scores to: Anger: 1, Surprise: 3, Shame: 9, Anticipation: 0.} \\
\\ 
\\ 
\\ 
\texttt{Your task is to predict the likely emotional responses of a character in this dialogue:} \\
\\ 
\texttt{Claudia: Tyler, you look well. Teaching suits you.} \\
\texttt{Tyler: Yeah, well, someone's gotta do the dirty work.} \\
\texttt{Claudia: I didn't mean it like that. I...} \\
\texttt{Tyler: I know what you meant. You always do, don't you?} \\
\texttt{Claudia: Not always, no. But I try to understand, at least.} \\
\texttt{Tyler: Understand? Like you understood when you left me alone?} \\
\texttt{Claudia: That's not fair, Ty. I had to...} \\
\texttt{Tyler: What? Chase your dreams? Leave your past behind?} \\
\texttt{Claudia: I never left you behind, Tyler. You were always...} \\
\texttt{Tyler: Always what? An afterthought? A burden?} \\
\texttt{Claudia: No, never a burden. You were my responsibility.} \\
\texttt{Tyler: Responsibility? Is that why you're here now? After all these years?} \\
\texttt{Claudia: I'm here because I...I wanted to see you.} \\
\texttt{Tyler: Or maybe you're here because you need something. Like always.} \\
\texttt{Claudia: That's not true, Tyler. I...} \\
\texttt{Tyler: Save it, Claudia. I've heard enough of your lies.} \\
\texttt{Claudia: I'm not lying, Tyler. I...} \\
\texttt{Tyler: You abandoned me, Claudia. You can't just waltz back into my life and pretend everything is okay.} \\
\texttt{Claudia: I'm not pretending, Tyler. I'm trying to make things right.} \\
\texttt{Tyler: Make things right? How? By throwing money at the problem?} \\
\texttt{Claudia: What are you talking about?} \\
\texttt{Tyler: Don't play dumb, Claudia. I know about the inheritance.} \\
\texttt{Claudia: I...I didn't know you knew.} \\
\texttt{Tyler: Of course you didn't. You never bothered to find out, did you?} \\
\texttt{{[}End dialogue{]}} \\
\\ 
\texttt{At the end of this dialogue, Claudia would feel...} \\
\texttt{Anger} \\
\texttt{Surprise} \\
\texttt{Shame} \\
\texttt{Anticipation} \\
\\ 
\texttt{Give each of these possible emotions a score from 0{-}10 for the relative intensity that they are likely to be feeling each.} \\
\\ 
\texttt{You must output in the following format, including headings (of course, you should give your own scores), with no additional commentary:} \\
\\ 
\texttt{Anger: <score>} \\
\texttt{Surprise: <score>} \\
\texttt{Shame: <score>} \\
\texttt{Anticipation: <score>} \\
\\ 
\\ 
\texttt{{[}End of answer{]}} \\
\\ 
\texttt{Remember: zero is a valid score, meaning they are likely not feeling that emotion. You must score at least one emotion > 0.} \\
\\ 
\texttt{Your answer:} \\
            }
        }
    }
     \\ \\

    \theutterance \stepcounter{utterance}  
    & \multicolumn{4}{p{0.6\linewidth}}{
        \cellcolor[rgb]{0.561,0.737,0.561}{
            \makecell[{{p{\linewidth}}}]{
                \texttt{\tiny{[P1$\rangle$GM]}}
                \texttt{Anger: 1} \\
\texttt{Surprise: 3} \\
\texttt{Shame: 9} \\
\texttt{Anticipation: 0} \\
            }
        }
    }
    & & \\ \\

    \theutterance \stepcounter{utterance}  
    & & & \multicolumn{2}{p{0.3\linewidth}}{
        \cellcolor[rgb]{0.9,0.9,0.9}{
            \makecell[{{p{\linewidth}}}]{
                \texttt{\tiny{[GM$|$GM]}}
                \texttt{Anger: 1} \\
\texttt{Surprise: 3} \\
\texttt{Shame: 9} \\
\texttt{Anticipation: 0} \\
            }
        }
    }
    & & \\ \\

    \theutterance \stepcounter{utterance}  
    & & & \multicolumn{2}{p{0.3\linewidth}}{
        \cellcolor[rgb]{0.9,0.9,0.9}{
            \makecell[{{p{\linewidth}}}]{
                \texttt{\tiny{[GM$|$GM]}}
                \texttt{Anger: 3} \\
\texttt{Surprise: 6} \\
\texttt{Shame: 8} \\
\texttt{Anticipation: 0} \\
            }
        }
    }
    & & \\ \\

    \theutterance \stepcounter{utterance}  
    & & & \multicolumn{2}{p{0.3\linewidth}}{
        \cellcolor[rgb]{0.9,0.9,0.9}{
            \makecell[{{p{\linewidth}}}]{
                \texttt{\tiny{[GM$|$GM]}}
                \texttt{game\_result = LOSE} \\
            }
        }
    }
    & & \\ \\

    \theutterance \stepcounter{utterance}  
    & & \multicolumn{4}{p{0.6\linewidth}}{
        \cellcolor[rgb]{0.95,0.95,0.95}{
            \makecell[{{p{\linewidth}}}]{
                \texttt{\tiny{[P1$\leftrightarrow$P]}}
                \texttt{Post{-}episode Reflexion (not sent to GM):} \\
\texttt{New plan: My previous attempt failed because I misinterpreted the "LOSE" outcome as a validation of my specific scores (Shame: 9, Surprise: 3), when in fact the system rejected the entire response. The dialogue reveals that Tyler has just exposed Claudia's secret regarding the inheritance, and her final line ("I...I didn't know you knew") indicates a moment of acute, fresh realization. While "Shame" is a strong underlying theme due to the accusations of abandonment, the immediate emotional spike at the very end of the dialogue is "Surprise" at the revelation that her secret is out. I incorrectly assumed the surprise had already subsided. Next time, I will prioritize the immediate reaction to the final piece of information revealed in the text. I will assign the highest score to "Surprise" to reflect the shock of the inheritance reveal, followed by "Shame" for the context of the accusations, while keeping "Anger" and "Anticipation" low. I will adjust the scores to: Anger: 1, Surprise: 9, Shame: 7, Anticipation: 0.} \\
            }
        }
    }
    & \\ \\

\end{supertabular}
}

\end{document}
